\section{Conclusion}
\label{sec:conclusion}

We turned a detector score into a statement about what no detector can achieve
on a given window, by measuring a two-sided indistinguishability bracket per
service pair on real QUIC wire data. At the coarsest pre-registered channel, 336
of 990 pairs (33.9\%) admit an interval whose upper side is finite-sample and
whose lower side is a classifier-based bound at a nominal level. The horizon is
short and its failure arithmetical rather than a property of the traffic: the
radius grows with the alphabet while the threshold does not, so the upper-side
set falls from 340 pairs at $M{=}1$ to none at $M{=}8$. Where it is lost, a
divisibility bound keeps the large-$k$ regime non-empty, and the missing-service
mass bounds what a bracket over observed services can claim.

Three directions follow: full-corpus confirmation of the occupancy pass, a
design clearing a known-truth gate while retaining the data, and other protocols
and channels, the frontier being a function of resolution, sample size, and
distance.
